\section{Gratian's Opening Distinctions}

Between Isidore and Aquinas stands the book that made the taxonomy \emph{canonical} in the strict sense: the \emph{Concordia discordantium canonum} of the Bologna master Gratian, compiled around 1140, called simply the \emph{Decretum}. It became the first volume of the \emph{Corpus iuris canonici}, the schoolbook on which every canonist to 1917 was trained; and it opens---before popes, before councils, before a single disciplinary canon---with the taxonomy of law. The Church's own lawbook begins by locating the Church's law inside a scheme larger than the Church's authority.\endnote{Gratian, \emph{Decretum}, first part, quoted in this section from the online transcription of Emil Friedberg's edition (\emph{Corpus iuris canonici}, editio Lipsiensis secunda, vol.~1, Leipzig 1879) published by the M\"unchener DigitalisierungsZentrum of the Bavarian State Library; the chapter pages quoted were fetched, hashed, and read 2026-07-25. Two ceilings govern everything below. First, Friedberg's is a nineteenth-century edition of the vulgate \emph{Decretum}; modern scholarship (Winroth and after) distinguishes recensions of the text, and nothing here depends on which stratum a given dictum belongs to. Second, the \emph{Decretum} was a school text that acquired immense authority by use; it was never promulgated as universal legislation, and it is cited here as the medium of reception, not as enacted law. Its Isidore is quoted as Gratian's Isidore, whose wording sometimes differs from Lindsay's text (e.g.\ \latin{Transire per agrum alienum} for Lindsay's \latin{Transire per alienum}); the differences are transmission history, not errors introduced by this article.}

\subsection{Distinction 1: two rules, and an equation}

The first words of the first distinction are Gratian's own dictum, and they say more than they seem to: \latin{Humanum genus duobus regitur, naturali uidelicet iure et moribus}---the human race is ruled by two things, namely by natural law and by usages. \latin{Ius naturae est, quod in lege et euangelio continetur, quo quisque iubetur alii facere, quod sibi uult fieri, et prohibetur alii inferre, quod sibi nolit fieri}---the law of nature is what is contained in the Law and the Gospel, by which each is commanded to do to another what he would have done to himself, and forbidden to inflict on another what he would not have inflicted on himself---sealed with Matthew 7:12: ``this is the law and the prophets.''\endnote{D.1 pr.\ (dictum Gratiani), followed immediately by the rubric introducing Isidore, \emph{Etymologies} V. All Gratian quotations preserve the Friedberg orthography (\latin{uidelicet}, \latin{euangelio}).}

Two moves happen in those two sentences. First, the whole of human governance is reduced to a pair: natural law and \latin{mores}---custom and enactment together, the entire apparatus of positive law, gathered under ``usages.'' Second, and more consequentially, the natural law is \emph{identified by its content with revelation's moral core}: what nature teaches is ``contained in the Law and the Gospel'' and summarized by the golden rule. Gratian then reproduces Isidore's chapters as his opening \emph{auctoritates}---c.~1 giving the divine/human division verbatim (\latin{Diuinae natura, humanae moribus constant \ldots\ Fas lex diuina est: ius lex humana}), c.~7 the definition of \latin{ius naturale}---and glosses them with a dictum that makes the compression explicit: from these words \latin{datur intelligi, in quo differant inter se lex diuina et humana, cum omne quod fas est, nomine diuinae uel naturalis legis accipiatur, nomine uero legis humanae mores iure conscripti et traditi intelligantur}---everything that is \latin{fas} is received under the name of \emph{divine or natural} law; under the name of human law are understood usages written down and handed on as law.\endnote{D.1 c.1 with the dictum post c.1; D.1 c.7. The chapter list of D.1 (cc.~1--12) reproduces Isidore's V.ii--ix almost seriatim: what is law, what custom, the species of \latin{ius}, natural, civil, of nations, military, public, Quirite.}

Note what Gratian has \emph{not} done: he has not distinguished natural law from divine positive law. \latin{Diuina uel naturalis lex} is, for the opening distinctions, one block---exactly Isidore's compression, now welded to Scripture. The scholastics would spend the next century taking that block apart (the canonists' glosses on this very page of the \emph{Decretum} enumerate multiple senses of \latin{ius naturale}); Aquinas's q.~91, with its separate articles for natural and divine law, is the mature repair. The 1983 Code, with its careful phrase \latin{legis divinae naturalis aut positivae} (c.~199 1°), is the repair become statute. A reader who wants to know why the Code must say ``divine law, natural \emph{or positive}'' is looking at the scar where Gratian's equation was surgically divided.

\subsection{Distinction 4: promulgation, and reception by use}

Distinction 4 carries the \emph{Decretum}'s doctrine of what makes an enactment a law. The dictum opening the distinction takes over Isidore's purpose-clause: laws are instituted \latin{ut} human audacity be checked and the capacity to harm restrained. Chapter 2 is Isidore's checklist (\latin{Erit autem lex honesta, iusta, possibilis \ldots}) with Gratian's pointed addendum: these things are to be weighed \emph{when law is being instituted}, for once laws have been instituted, \latin{non erit liberum iudicare de ipsis, sed oportebit iudicare secundum ipsas}---it will not be free to judge \emph{about} them, but one will have to judge \emph{according to} them.\endnote{D.4 pr.; D.4 c.2 with the dictum post c.2 (quoting Augustine, \emph{De vera religione}, as c.3). Gratian's addendum is an early and exact statement of the difference between legislative deliberation and adjudication under law---a rule-of-law principle inside a twelfth-century church collection.}

Then, after Augustine's chapter making the same point, the dictum the canonical tradition never forgot: \latin{Leges instituuntur, cum promulgantur, firmantur, cum moribus utentium approbantur}---laws are \emph{instituted} when they are promulgated; they are \emph{confirmed} when they are approved by the usages of those who use them. As some laws have today been abrogated by the contrary usage of users, so laws are confirmed by users' usage; whence the decree of Pope Telesphorus that clergy fast from Quinquagesima, never approved by use, does not convict those who do otherwise of transgression.\endnote{D.4 dictum post c.3, with the Telesphorus example. The modern Code's promulgation rule (c.~7: \latin{Lex instituitur cum promulgatur}, a verbatim descendant of this dictum's first clause) keeps the institution half; the confirmation-by-use half survives, transformed and disciplined, in the custom canons (cc.~23--28) and in the doctrine of desuetude---no longer as a general condition of validity. The difference between Gratian's world and the Code's on exactly this point is real, and section 9 treats it.}

This is the taxonomy's pastoral realism: positive law lives between two waters, authority above and reception below. Gratian will not let either swallow the other---promulgation institutes without use, but use can unmake what promulgation made.

\subsection{Distinctions 8 and 9: the hierarchy stated as law}

Distinction 8 opens by dividing what D.1 had joined: \latin{Differt etiam ius naturae a consuetudine et constitutione}---the law of nature differs from custom and enactment. By the law of nature all things are common to all---as in the believers of Acts 4, and as the philosophers handed down (Plato's city, ``in which no one knows his own attachments''); by the law of custom or enactment, \emph{this is mine, that another's}.\endnote{D.8 pr. Property, on this classical view, is an institution of human positive law which natural law permits but does not itself institute---Isidore's ``common possession of all things'' finally cashed out. Augustine's chapter follows: by what right do you defend the Church's villas, divine or human? Human right, the right of emperors and kings---\latin{Tolle iura imperatoris, et quis audet dicere: hec uilla mea est?} (D.8 c.1).}

And then Gratian states the hierarchy as flatly as it has ever been stated: \latin{Dignitate uero ius naturale simpliciter preualet consuetudini et constitutioni. Quecunque enim uel moribus recepta sunt, uel scriptis comprehensa, si naturali iuri fuerint aduersa, uana et irrita sunt habenda}---in dignity the natural law simply prevails over custom and enactment: whatever has been received in usages or comprehended in writings, if it is adverse to natural law, is to be held vain and void.\endnote{D.8, dictum Gratiani opening the second part (after c.1). The distinction's remaining chapters (Cyprian, Augustine, Gregory) hammer the corollary against custom: ``the Lord said \emph{I am the truth}; he did not say \emph{I am custom}''; custom without truth is \latin{uetustas erroris}, the old age of error (D.8 cc.~4--9), closing with the dictum: \latin{Liquido igitur apparet, quod consuetudo naturali iuri postponitur}---custom is subordinate to natural law.}

Distinction 9 does the same for enactment: \latin{Quod autem constitutio naturali iuri cedat multiplici auctoritate probatur}---that enactment yields to natural law is proved by manifold authority---and, after eleven chapters of Augustine, concludes: since by the natural law nothing is commanded but what God wills, and nothing forbidden but what God prohibits, and since the divine laws stand by nature, therefore whatever is shown contrary to the divine will, or to canonical Scripture, or to the divine laws, is also adverse to the natural law; \latin{Constitutiones ergo uel ecclesiasticae uel seculares, si naturali iuri contrariae probantur, penitus sunt excludendae}---enactments, whether \emph{ecclesiastical} or secular, if they are proved contrary to natural law, are to be entirely excluded.\endnote{D.9 pr.\ and dictum post c.11. The inclusion of \latin{ecclesiasticae} deserves its emphasis: the first pages of the Church's own classical lawbook apply the supremacy of natural law to the Church's enactments by name, not only to the state's. The verb is also worth exactness: \latin{excludendae}---to be excluded, set aside as without force---a juridical conclusion about validity, in line with D.8's \latin{uana et irrita}, not a general theory of resistance.}

The architecture of the modern canons is now fully prefigured. That no custom contrary to divine law can obtain the force of law (c.~24 §1); that civil laws are received \latin{quatenus iuri divino non sint contrariae} (cc.~22, 1290); that dispensation reaches merely ecclesiastical law and stops there (c.~85); that the canons which repeat the old law are to be weighed \latin{ratione etiam canonicae traditionis habita}---with account taken of the canonical tradition (c.~6 §2)\endnote{CIC c.~6 §2: \latin{Canones huius Codicis, quatenus ius vetus referunt, aestimandi sunt ratione etiam canonicae traditionis habita.} Latin and English deliveries fetched, hashed, and read 2026-07-25. The canon is the Code's own instruction to read it the way this article is reading it: continuity with the tradition of which Gratian is the classical carrier is an interpretive norm of the current law, not an antiquarian preference.}---all of it is Gratian's opening distinctions, restated with a codifier's economy. What the \emph{Decretum} argued from authorities, the Code enacts as rules; and it is to the Code that we now turn.
